% !PRINTED 172
\weberpage{172}

If we therefore replace $\omega$ by $2\omega$ in (17) and extract the square root, it follows that:

\weq{19}{q = \frac{1}{4}\frac{\varkappa_1}{\varkappa_1'}\, e^{-\frac{1}{2}\frac{G(\varkappa_1{}^2)}{F(\varkappa_1{}^2)}}.}

Here the domain of convergence now extends over the unit circle lying in the plane $\varkappa_1{}^2$, and this circle corresponds to the whole plane $\varkappa^2$. For the real part of $\varkappa'$ is positive in the \gesperrt{whole plane} of the $\varkappa^2$, with the exception of the two banks of the cut running from $+1$ to $+\infty$, on which $\varkappa'$ is purely imaginary. From this it follows that the absolute value of $\varkappa_1$ is less than 1, and that the periphery of the unit circle in the plane $\varkappa_1$ corresponds to the two banks of this cut.

But one can improve the convergence of these expressions still further by a renewed application of the \gesperrt{Landen} transformation. If we denote the result of a further doubling of $\omega$ by the index 2, then it follows that:

\weq{20}{\begin{aligned}
\sqrt{\varkappa_2} &= \frac{1-\sqrt{\varkappa'}}{1+\sqrt{\varkappa'}}, \\
4K_2 &= (1+\sqrt{\varkappa'})^2 K, \quad K_2' = (1+\sqrt{\varkappa'})^2 K', \\
q &= \sqrt{\frac{\varkappa_2}{4\varkappa_2'}}\, e^{-\frac{1}{4}\frac{G(\varkappa_2{}^2)}{F(\varkappa_2{}^2)}}.
\end{aligned}}

in which the domain of convergence now covers the plane of the $\varkappa^2$ \gesperrt{twice} (with a branching at the point $\varkappa^2=1$).

This operation can be continued, by deriving from $\varkappa_2$, according to formula (17), a new quantity $\varkappa_3$ through doubling of $\omega$, and from this in the same way $\varkappa_4$, and so on. One then obtains a sequence of expressions for $q$ whose convergence becomes ever better. For an arbitrary $\nu$ we have:

\weq{21}{q = \sqrt[2^{\nu-1}]{\frac{\varkappa_\nu}{4\varkappa_\nu'}}\, e^{-\frac{G(\varkappa_\nu{}^2)}{2^\nu F(\varkappa_\nu{}^2)}}.}

The first equation (18)

\[
\varkappa_1 = \frac{\varkappa^2}{(1+\varkappa')^2}
\]

shows that, so long as the absolute value of $\varkappa^2$ is less than 1, and consequently the absolute value of $1+\varkappa'$ greater than 1, the absolute value of $\varkappa_1$ is less than that of $\varkappa^2$, and consequently
%CONT>>
